\chapter{高頻度取引とマーケットメイクの研究}
\label{ch:hft}

第 \ref{ch:mm} 章で、板の仕組み、逆選択と在庫のモデルを説明した。この章では、高頻度の取引会社が実際に何をし、どれだけ儲け、市場にどんな影響を与えているかの実証研究をまとめる。

\section{新しいマーケットメイカー}

\begin{paperbox}[Menkveld（2013）「高頻度取引と新しいマーケットメイカー」\cite{menkveld2013}]
欧州の株式で、既存の取引所と、急成長していた新しい取引所（Chi-X）の両方で売買していた一つの大手の高頻度の取引会社の全取引を分析した。この会社は、既存の取引所の取引の 8.1\%、新しい取引所の取引の 64.4\% に参加し、どちらでも取引の約 5 分の 4 は、自分の指値が相手の注文で成立する「受け身の」取引（流動性の供給）だった。在庫の値動きでは損をし、スプレッドで稼いでいた。シャープレシオの計算は、資本のコストをどう仮定するかに強く依存する、と要旨は述べている。
\end{paperbox}

Deep Research の答え 1 は、この論文の値として「1 銘柄・1 日に 1,397 回の取引、1 回 0.88 ユーロの利益、シャープレシオ 9.35」を挙げたが、要旨には出てこず、本文で確かめられなかったので、この報告書では使わない。

この研究は、高頻度取引の大きな部分が、昔の取引所の立会場のマーケットメイカーの役割を、計算機で引き継いだものであることを示した。

\section{流動性と価格の発見}

\begin{itemize}
\item \textbf{Hendershott・Jones・Menkveld（2011）}\cite{hjm2011}：ニューヨーク証券取引所が 2003 年に気配の自動配信を導入し、アルゴリズム取引が増えた変化を使って、アルゴリズム取引がスプレッドを狭め、気配の情報を増やし、流動性を改善したことを示した。
\item \textbf{Brogaard・Hendershott・Riordan（2014）}\cite{bhr2014}：米国株の高頻度の取引会社の売買を分析し、高頻度の取引会社は、全体として、価格の恒久的な変化の方向に売買し、一時的な価格の歪みの反対に売買することで、価格の発見を助けていることを示した。一方で、高頻度の取引会社が流動性を供給する取引は、逆選択を受ける（相手に情報がある）傾向があった。
\end{itemize}

\section{フラッシュ・クラッシュ}

\begin{paperbox}[Kirilenko・Kyle・Samadi・Tuzun（2017）\cite{kkst2017}]
2010 年 5 月 6 日、米国の株価指数の先物（E-mini \SP{}）が数分で急落し、すぐに戻った。先物の取引の口座ごとのデータを分析し、高頻度の取引会社は急落の引き金ではなかったが、大口の売りの注文が市場に入ったとき、在庫を素早く互いに売り渡し（「熱いじゃがいも」）、売りの圧力を増幅したことを示した。高頻度の取引会社の売買のやり方は、急落の日も普段と変わらなかった。
\end{paperbox}

\section{速さの競争}

\begin{paperbox}[Budish・Cramton・Shim（2015）\cite{bcs2015}]
株価指数の先物と、その ETF の価格は、秒の単位ではほぼ完全に連動するが、ミリ秒の単位では連動が崩れ、機械的な裁定の機会が生じる。この機会は速い者だけが取れるので、速さへの投資の競争が起きるが、機会の大きさは速さの競争で小さくならない。連続的に売買を成立させる市場の仕組みの欠陥であり、一定の間隔でまとめて成立させる方式（頻繁なバッチ・オークション）が解決策だと論じた。
\end{paperbox}

\begin{paperbox}[Aquilina・Budish・O'Neill（2022）\cite{abo2022}]
ロンドン証券取引所のすべての注文の記録（成立しなかった注文や、失敗した注文を含む）から、同じ価格の変化に向けて複数の会社が同時に注文を送る「速さの競争」を数えた。FTSE100 の銘柄では、こうした競争が頻繁に起き、取引量のかなりの部分を占める。遅延の裁定は、投資家にとっての取引のコストを押し上げる「税」のように働き、それをなくすと流動性のコストが下がると推計した。
\end{paperbox}

\section{利益の集中}

Baron・Brogaard・Hagströmer・Kirilenko（2019）\cite{bbhk2019} は、スウェーデンの株式市場の取引会社のデータで、高頻度取引の利益が少数の会社に集中し、速さの順位が利益を強く左右し、新しく参入した会社は利益が小さく退出しやすいことを示した。高頻度取引は、規模と速さの経済が強く働く、勝者が多くを取る市場である。

\section{板の情報による短期の予測}

Cont・Kukanov・Stoikov（2014）\cite{cks2014} は、最良気配の注文の量の増減の偏り（注文の流れの不均衡）が、数秒〜数分の価格の変化を線形によく説明し、その傾きは板の厚さに反比例することを示した。短期の価格の予測とインパクトのモデルの基本になっている。

\section{確立していること}

\begin{keybox}[まとめ]
\begin{itemize}
\item 高頻度取引の大きな部分はマーケットメイクであり、平常時には、スプレッドを狭め、価格の発見を助けている。
\item 収益は、極めて多くの小さな取引の積み重ねで、シャープレシオは非常に高いが、容量は小さく、少数の速い会社に集中する。
\item 市場が混乱すると、流動性の供給を引っ込め、急落を増幅することがある。
\item 速さの競争は、市場の仕組み（連続的な売買の成立）が生んでいて、社会的な価値の小さい投資を招いている、という議論がある。
\end{itemize}
\end{keybox}

個人や大学の研究がこの分野で競うことは事実上できないが、決算発表の直後の数秒〜数分は、この人たちが情報を価格に入れる時間であることを知っておく価値はある（第 \ref{ch:mm} 章の「この研究との関係」）。
